\section{Decisiones de Diseño y Arquitectura}

Para que la simulación de este sistema de infusión sea realista y permita evaluar correctamente las condiciones de falla requeridas (como el desvío del 10\% del caudal), se tomaron decisiones de diseño fundamentales que extienden la propuesta básica:

\subsection{Arquitectura de Lazo Cerrado (Closed-Loop)}
El flujo de datos se estructuró separando claramente la orden ideal de la realidad física:
\begin{itemize}
    \item El \textbf{Controlador} no asume que la bomba funciona perfecto. Le envía el \texttt{caudalObjetivo} al \textbf{Actuador}.
    \item El \textbf{Actuador} procesa esa orden aplicándole un retraso mecánico y establece el \texttt{caudalReal} físico.
    \item El \textbf{Sensor de Flujo} lee ese \texttt{caudalReal}, le aplica un ruido o error de medición esperado en cualquier instrumento, y le devuelve el \texttt{caudalMedido} al \textbf{Controlador}.
\end{itemize}
Gracias a esta realimentación, el Controlador puede comparar la orden original con lo que el sensor está midiendo, y así detectar discrepancias para disparar las alarmas media o crítica.

\subsection{Modelado Estocástico de Eventos}
Se decidió modelar los eventos externos e impredecibles utilizando distribuciones de probabilidad teóricas:
\begin{itemize}
    \item \textbf{Llegada de eventos (Órdenes y Confirmaciones):} Se utiliza una distribución \textbf{Exponencial} para modelar el tiempo entre la llegada de nuevas prescripciones médicas ($G_{om}$) y la intervención del personal de enfermería ($G_{ce}$). Esto asume un Proceso de Poisson, ideal para eventos independientes sin memoria.
    \item \textbf{Caudal y Latencias:} Se utiliza una distribución \textbf{Uniforme} para el caudal objetivo ($[0, 200]$ ml/h) y para la latencia del actuador ($[0, 0.5]$ s). Esto garantiza que el sistema sea evaluado en todos sus extremos posibles sin sesgar las pruebas hacia un valor particular.
    \item \textbf{Ruido del Sensor:} Se utiliza una distribución \textbf{Normal} (con una desviación estándar del 20\%) para alterar la lectura del caudal real. Según el Teorema del Límite Central, la suma de pequeñas perturbaciones mecánicas y de sensado tiende a una distribución Gaussiana, lo que forzará al controlador a gestionar desvíos realistas.
\end{itemize}

\section{Modelos Atómicos Definidos}

\subsection{Generador de Órdenes Médicas ($G_{om}$)}
Este componente modela la llegada estocástica de nuevas prescripciones médicas que alteran el caudal objetivo de la bomba.
Matemáticamente, el modelo se define como:
\[ G_{om} = \langle X_{G_{om}}, Y_{G_{om}}, S_{G_{om}}, \delta_{int}, \delta_{ext}, \lambda, ta \rangle \]
donde:
\begin{itemize}
    \item $X_{G_{om}} = \emptyset$ (No posee puertos de entrada).
    \item $Y_{G_{om}} = [0, 200] \times \{out\_caudal\_obj\}$
    \item $S_{G_{om}} = [0, 200] \times \mathbb{R}_{\geq 0}$, donde $s = (caudalObjetivo, \sigma)$.
    \item $s_0 = (\text{Uniforme}(0, 200), \text{Exponencial}(1 / 300.0))$.
    \item $\delta_{ext}(s, e, x) = \emptyset$
    \item \(\delta_{int}(caudalObjetivo, \sigma) = (\text{Uniforme}(0, 200), \text{Exponencial}(1 / 300.0))\).
    \item \(\lambda(caudalObjetivo, \sigma) = (caudalObjetivo, out\_caudal\_obj)\).
    \item \(ta(caudalObjetivo, \sigma) = \sigma\).
\end{itemize}

\subsection{Actuador de la Bomba ($A$)}
Este componente modela la respuesta física y mecánica de la bomba. Presenta una latencia inherente al motor.
Matemáticamente, el modelo se define como:
\[ A = \langle X_A, Y_A, S_A, \delta_{int}, \delta_{ext}, \lambda, ta \rangle \]
donde:
\begin{itemize}
    \item $X_A = ([0, 200] \times \{in\_caudal\_obj\}) \cup (\{detenerBomba\} \times \{in\_stop\})$
    \item $Y_A = [0, 200] \times \{out\_caudal\_real\}$
    \item $S_A = [0, 200] \times \mathbb{R}_{\geq 0} \cup \{\infty\}$, donde $s = (caudalReal, \sigma)$.
    \item $s_0 = (0.0, \infty)$.
    \item Transición Externa:
    \[
    \delta_{ext}((caudalReal, \sigma), e, (Xv, port)) = 
    \begin{cases} 
    (Xv, \text{Uniforme}(0, 0.5)) & \text{si } port = in\_caudal\_obj \\
    (0.0, \text{Uniforme}(0, 0.5)) & \text{si } port = in\_stop 
    \end{cases}
    \]
    \item $\delta_{int}(caudalReal, \sigma) = (caudalReal, \infty)$
    \item $\lambda(caudalReal, \sigma) = (caudalReal, out\_caudal\_real)$.
    \item $ta(caudalReal, \sigma) = \sigma$.
\end{itemize}

\subsection{Sensor de Flujo ($G_{sf}$)}
Este modelo representa el dispositivo encargado de medir el caudal real. Toma muestras periódicas inyectando un ruido Gaussiano del 20\%.
Matemáticamente, el modelo se define como:
\[ G_{sf} = \langle X_{G_{sf}}, Y_{G_{sf}}, S_{G_{sf}}, \delta_{int}, \delta_{ext}, \lambda, ta \rangle \]
donde:
\begin{itemize}
    \item $X_{G_{sf}} = [0, 200] \times \{in\_caudal\_real\}$
    \item $Y_{G_{sf}} = [0, 200] \times \{out\_caudal\_medido\}$
    \item $S_{G_{sf}} = [0, 200] \times [0, 200] \times \mathbb{R}_{\geq 0} \cup \{\infty\}$, donde $s = (caudalReal, caudalMedido, \sigma)$.
    \item $s_0 = (0.0, 0.0, \infty)$.
    \item Transición Externa (dividida para evitar desbordamiento):
    \[
    \delta_{ext}((caudalReal, caudalMedido, \sigma), e, (Xv, port)) = 
    \]
    \[
    \begin{cases} 
    (Xv, \text{Normal}(Xv, (0.20 \cdot Xv)^2), \sigma - e) \\ 
    \quad \text{si } port = in\_caudal\_real \land \sigma < \infty \\[2ex]
    (Xv, \text{Normal}(Xv, (0.20 \cdot Xv)^2), 0.0) \\ 
    \quad \text{si } port = in\_caudal\_real \land \sigma = \infty
    \end{cases}
    \]
    \item $\delta_{int}(caudalReal, caudalMedido, \sigma) = (caudalReal, \text{Normal}(caudalReal, (0.20 \cdot caudalReal)^2), 1.0)$
    \item $\lambda(caudalReal, caudalMedido, \sigma) = (caudalMedido, out\_caudal\_medido)$.
    \item $ta(caudalReal, caudalMedido, \sigma) = \sigma$.
\end{itemize}

\subsection{Generador de Confirmaciones del Enfermero ($G_{ce}$)}
Componente que simula los retrasos operativos humanos, enviando la confirmación de la visualización de una alarma.
Matemáticamente, el modelo se define como:
\[ G_{ce} = \langle X_{G_{ce}}, Y_{G_{ce}}, S_{G_{ce}}, \delta_{int}, \delta_{ext}, \lambda, ta \rangle \]
donde:
\begin{itemize}
    \item $X_{G_{ce}} = \emptyset$
    \item $Y_{G_{ce}} = \{confirmacionEnfermero\} \times \{out\_conf\}$
    \item $S_{G_{ce}} = \mathbb{R}_{\geq 0}$, donde $s = \sigma$.
    \item $s_0 = \text{Exponencial}(1 / 8.0)$.
    \item $\delta_{ext}(s, e, x) = \emptyset$.
    \item $\delta_{int}(\sigma) = \text{Exponencial}(1 / 8.0)$.
    \item $\lambda(\sigma) = (confirmacionEnfermero, out\_conf)$.
    \item $ta(\sigma) = \sigma$.
\end{itemize}

\subsection{Generador de Fin de Bolsa ($G_{fb}$)}
Este componente modela la bolsa de infusión. Rastrea el volumen restante de líquido
en base al caudal medido por el Sensor de Flujo y emite una señal de alerta cuando
estima que quedan 60 segundos o menos para agotar la bolsa.
Se asume una capacidad de bolsa $V_0 = 500$ ml y un umbral de alerta $T_a = 60$ s.

Matemáticamente, el modelo se define como:
\[ G_{fb} = \langle X_{G_{fb}}, Y_{G_{fb}}, S_{G_{fb}}, \delta_{int}, \delta_{ext}, \lambda, ta \rangle \]
donde:
\begin{itemize}
    \item $X_{G_{fb}} = ([0, 200] \times \{in\_caudal\_medido\}) \cup (\{\top\} \times \{in\_rellenar\})$
    \item $Y_{G_{fb}} = \{\top\} \times \{out\_fin\_bolsa\}$
    \item $S_{G_{fb}} = \{\text{MONITOREANDO}, \text{ESPERANDO\_RELLENO}\}
           \times \mathbb{R}_{\geq 0}
           \times [0, 200]
           \times \mathbb{R}_{\geq 0} \cup \{\infty\}$,
           donde $s = (fase, volRest, q, \sigma)$.

    Las variables del estado representan:
    \begin{itemize}
        \item $fase$: fase actual del componente.
        \item $volRest$: volumen restante de líquido en la bolsa (ml).
        \item $q$: último caudal medido informado por el Sensor de Flujo (ml/h).
        \item $\sigma$: tiempo hasta la próxima transición interna.
    \end{itemize}

    \item $s_0 = (\text{MONITOREANDO},\; V_0,\; 0.0,\; \infty)$.

    \item Funciones auxiliares. $vol'$ descuenta el volumen consumido durante $e$ segundos al caudal $q$, y $\hat{\sigma}$ calcula cuánto falta para emitir la alerta dado un volumen y un caudal:
    \[
    vol'(volRest, q, e) = \max\!\left(volRest - q \cdot \frac{e}{3600},\; 0\right)
    \]
    \[
    \hat{\sigma}(v, q) =
    \begin{cases}
    \infty & \text{si } q = 0 \\
    0 & \text{si } \dfrac{v}{q / 3600} \leq T_a \\[1ex]
    \dfrac{v}{q / 3600} - T_a & \text{en otro caso}
    \end{cases}
    \]

    \item Transición Externa. Sea $v' = vol'(volRest, q, e)$:
    \[
    \delta_{ext}((fase, volRest, q, \sigma), e, (Xv, port)) =
    \]
    \[
    \begin{cases}
    (\text{ESPERANDO\_RELLENO},\; v',\; Xv,\; \infty) \\
    \quad \text{si } port = in\_caudal\_medido \land fase = \text{ESPERANDO\_RELLENO} \\[2ex]
    (\text{MONITOREANDO},\; v',\; Xv,\; \hat{\sigma}(v', Xv)) \\
    \quad \text{si } port = in\_caudal\_medido \land fase = \text{MONITOREANDO} \\[2ex]
    (\text{MONITOREANDO},\; V_0,\; q,\; \hat{\sigma}(V_0, q)) \\
    \quad \text{si } port = in\_rellenar
    \end{cases}
    \]

    \item $\delta_{int}(fase, volRest, q, \sigma)
           = (\text{ESPERANDO\_RELLENO},\; volRest,\; q,\; \infty)$

    \item Función de salida:
    \[
    \lambda(fase, volRest, q, \sigma) =
    \begin{cases}
    (\top,\; out\_fin\_bolsa) & \text{si } fase = \text{MONITOREANDO} \\
    \emptyset & \text{en otro caso}
    \end{cases}
    \]

    \item $ta(fase, volRest, q, \sigma) = \sigma$.
\end{itemize}

\subsection{Módulo de Alarmas ($M_a$)}
\textit{La especificación formal del procesamiento de prioridades de alertas y su temporización se desarrollará en la siguiente etapa del proyecto.}

\subsection{Controlador de Bomba ($C$)}
\textit{La especificación formal del núcleo de control lógico y de toma de decisiones del sistema se incorporará en la siguiente etapa del proyecto.}

\section{Especificación del Modelo Acoplado Global}
A continuación se detalla la infraestructura de interconexiones que da forma al sistema de lazo cerrado:

\begin{figure}[h]
    \centering
    % \includegraphics[width=\textwidth]{diagrama_acoplado.png}
    \textit{[Diagrama del modelo acoplado pendiente de inclusión.]}
    \caption{Diagrama del modelo acoplado global del sistema de infusión.}
    \label{fig:diagrama_acoplado}
\end{figure}

\begin{itemize}
    \item \textbf{Generador de Fin Bolsa ($G_{fb}$)} $\rightarrow$ Envía la señal \texttt{finBolsa} al \textbf{Controlador de Bomba}.
    \item \textbf{Generador de Confirmación ($G_{ce}$)} $\rightarrow$ Conecta \texttt{out\_conf} al \textbf{Controlador de Bomba}.
    \item \textbf{Generador de Órdenes ($G_{om}$)} $\rightarrow$ Conecta \texttt{out\_caudal\_obj} al \textbf{Controlador de Bomba}.
    \item \textbf{Controlador de Bomba} $\rightarrow$ Envía la consigna por su puerto hacia \texttt{in\_caudal\_obj} del \textbf{Actuador de la Bomba}, y \texttt{detenerBomba} hacia \texttt{in\_stop}. Despacha las alertas al \textbf{Módulo de Alarmas}.
    \item \textbf{Actuador de la Bomba ($A$)} $\rightarrow$ Comunica el \texttt{caudalReal} a través de \texttt{out\_caudal\_real} hacia \texttt{in\_caudal\_real} del \textbf{Sensor de Flujo}.
    \item \textbf{Sensor de Flujo ($G_{sf}$)} $\rightarrow$ Envía el \texttt{caudalMedido} (con ruido) desde \texttt{out\_caudal\_medido} hacia el \textbf{Controlador de Bomba} y el \textbf{Generador de Fin de Bolsa ($G_{fb}$)}.
\end{itemize}
